\DHead{III. 核心原则：行动框架}
\DLayer{第一层：筑牢智能时代基石}
\DPrinciple{原则一：坚持开放互联}{“万物智联，普惠共生”}
打破算法壁垒与语言障碍，推动多模态模型的跨文化交互。鼓励支持多语言的本地化大模型开发，坚持技术发展的非歧视性。让发展中国家能够平等接入全球 AI 资源，使技术连接成为缩小数字鸿沟的桥梁，而非加剧不平等的工具。

\DPrinciple{原则二：坚持安全可信}{“透明算法，信任基石”}
反对互联网碎片化与技术封闭。支持开源作为验证算法安全、检测数据偏见的首选试验田。建立可信、弹性、可持续的 AI 供应链安全机制，确保关键基础设施的中立性与稳定性，为全球智能社会的运行筑牢底线。

\DLayer{第二层：夯实伦理治理支柱}
\DPrinciple{原则三：坚持数据伦理}{“算法向善，责任共担”}
坚持数据隐私保护与算法公平性，防范技术滥用。推动 AI 治理遵循包容、敏捷、负责任的原则。鼓励开源社区探索可解释性 AI（XAI）的最佳实践，确保每一项 AI 决策都可追溯、可审计，始终服务于人类尊严与权利。

\DPrinciple{原则四：坚持多元包容}{“众智共创，价值共享”}
尊重文化多样性与认知能力的差异，反对任何形式的“算法歧视”。特别关注弱势群体在 AI 时代的权益，通过开源项目赋能边缘群体。支持多元化的开发者社群，让每一个声音都能在智能技术的演进中被听见、被尊重。

\DPrinciple{原则五：坚持全球协作}{“跨界协同，智能互联”}
视开源为连接全球智慧的网络。支持跨国界、跨机构的联合研发，特别是在应对气候变化、生物医药等全球性挑战领域。促进知识、算力与高质量数据集的自由流动，强化南南合作与多边主义，构建去中心化但互联互通的创新生态。

\DLayer{第三层：共创社会价值未来}
\DPrinciple{原则六：坚持绿色智能}{“低碳驱动，创新致远”}
优先支持应对气候变化、环境保护、清洁能源等议题的开源 AI 项目。倡导高效的模型训练范式与算力调度优化，降低 AI 运行的碳足迹。推动开源技术作为数字公共品的核心实现路径，在能源管理、环境监测中的广泛应用，确保智能技术的发展不以牺牲生态环境为代价。
